\documentclass[11pt]{article}
\usepackage[margin=1in]{geometry}
\usepackage{amsmath,amssymb,amsthm,mathtools}
\usepackage{graphicx}
\usepackage[hidelinks]{hyperref}
\usepackage{xcolor}

\newtheorem{theorem}{Theorem}
\newtheorem{lemma}[theorem]{Lemma}
\newcommand{\R}{\mathbb R}
\newcommand{\C}{\mathbb C}
\newcommand{\Pp}{\mathbb P}
\newcommand{\E}{\mathbb E}
\newcommand{\Conf}{\operatorname{Conf}}

\title{Centered Giambelli Compatibility Does Not\\
Characterize Determinantal Point Processes}
\author{Automatic mathematics run \texttt{fa\_banach\_001}}
\date{11 August 2026}

\begin{document}
\setlength{\emergencystretch}{3em}
\maketitle

\begin{center}
\fbox{\parbox{0.91\textwidth}{
\textbf{Status.} Candidate counterexample, likely valid, for expert review.
For every $R>0$, the simple point process consisting of two independent
Cauchy$(0,R)$ points satisfies every hypothesis of Conjecture~1 in
Bufetov--Lazag but is not determinantal.  Thus the conjecture as stated has a
full negative answer.
}}
\end{center}

\section{Source conjecture}

The source is A.~I. Bufetov and P. Lazag,
\emph{A determinantal point process governed by an integrable projection
kernel is Giambelli compatible}, arXiv:2111.05606 \cite{source}.
For a point configuration $X\subset\R$ and $R>0$, its centered
specialization is prescribed on the Newton power sums by
\begin{equation}\label{eq:rho}
 \rho^R(p_k)(X)=
 \sum_{x\in X}\frac1{(x+iR)^k}
 -\E_{\Pp}\sum_{x\in X}\frac1{(x+iR)^k},
 \qquad k\ge1.
\end{equation}
On printed page 5, Conjecture~1 asks whether every point process for which
$\rho^R$ is an $L^1$-specialization and $(\Pp,\rho^R)$ is Giambelli compatible
must be determinantal.  The source statement is reproduced below.

\begin{center}
\includegraphics[width=0.98\textwidth]{figures/open_problem_crop.png}
\end{center}

Here Giambelli compatibility means that, for every partition
$\lambda=(p_1,\ldots,p_d\mid q_1,\ldots,q_d)$ in Frobenius notation,
\begin{equation}\label{eq:giambelli}
 \E_{\Pp}[\rho^R(s_\lambda)]
 =\det\bigl(\E_{\Pp}[\rho^R(s_{(p_i\mid q_j)})]\bigr)_{i,j=1}^d.
\end{equation}

\section{Proof intuition}

The centering in \eqref{eq:rho} creates a degeneracy on harmonic measure.
The M\"obius map $x\mapsto (x+iR)^{-1}$ sends Cauchy$(0,R)$ measure to
uniform arclength on a circle.  Holomorphic mean value then replaces each
random transformed point by the center of that circle.  With two independent
points, centering subtracts exactly the power sums of those two center points.
Consequently, the expectation of the entire centered symmetric-function
specialization is evaluation at the zero alphabet.  Every nonempty Schur
expectation vanishes, so all Giambelli identities hold for a reason unrelated
to fermionic or determinantal structure.

To rule out a hidden determinantal representation, use only the first four
correlation functions.  A fixed-two iid process has correlation densities
$2,2,0,0$.  Were these principal minors of a kernel, the two- and three-minor
equations would force a normalized four-point kernel matrix into one of two
tournament forms.  The only form allowed by the three-minors has determinant
$-4$, whereas the four-point correlation requires determinant zero.

\section{Counterexample}

\begin{theorem}[Iid Cauchy counterexample]\label{thm:main}
Fix $R>0$, and let $\nu_R$ be the Cauchy probability measure
\begin{equation}\label{eq:cauchy}
 d\nu_R(x)=\frac{R}{\pi(x^2+R^2)}\,dx
 \qquad(x\in\R).
\end{equation}
Let $Y_1,Y_2$ be independent with law $\nu_R$, and let $\Pp$ be the law on
$\Conf(\R)$ of the random set $X=\{Y_1,Y_2\}$.  Then:
\begin{enumerate}
 \item $X$ is a simple two-point configuration almost surely;
 \item the specialization $\rho^R$ in \eqref{eq:rho} is an
 $L^1(\Conf(\R),\Pp)$-specialization;
 \item $(\Pp,\rho^R)$ is Giambelli compatible;
 \item $\Pp$ is not a determinantal point process, even if complex-valued,
 nonsymmetric measurable kernels and arbitrary dominating reference measures
 are allowed.
\end{enumerate}
Hence Conjecture~1 of \cite{source} is false.
\end{theorem}

\section{Proof}

Because $\nu_R$ is atomless, $Y_1\ne Y_2$ almost surely.  Thus the unordered
set $X$ is indeed a simple point configuration with exactly two points.

\subsection*{Cauchy measure becomes uniform circle measure}

Put
\begin{equation}\label{eq:transform}
 T(x)=\frac1{x+iR},\qquad
 c=-\frac{i}{2R},\qquad r=\frac1{2R}.
\end{equation}
If $\Theta$ is uniform on $(0,2\pi)$, then
$R\cot(\Theta/2)$ has law $\nu_R$, and a direct calculation gives
\begin{equation}\label{eq:circleparam}
 T\bigl(R\cot(\Theta/2)\bigr)
 =-\frac{i}{2R}+\frac{i e^{-i\Theta}}{2R}.
\end{equation}
Thus $T_\#\nu_R$ is uniform arclength on the circle
$\{c+r\zeta:|\zeta|=1\}$.  In particular, the one-variable mean-value
identity gives
\begin{equation}\label{eq:moments}
 \int_\R T(x)^k\,d\nu_R(x)=c^k,
 \qquad k\ge1.
\end{equation}

Let $t_j=T(Y_j)$.  More generally, for every polynomial
$F\in\C[z_1,z_2]$, independence and two applications of the mean-value
identity give
\begin{equation}\label{eq:twomean}
 \E F(t_1,t_2)=F(c,c).
\end{equation}

\subsection*{$L^1$ specialization and Giambelli compatibility}

Since $|t_j|\le R^{-1}$, \eqref{eq:moments} yields
\begin{equation}\label{eq:centeredpk}
 \rho^R(p_k)=t_1^k+t_2^k-2c^k,
 \qquad k\ge1.
\end{equation}
The algebra $\Lambda$ of symmetric functions is the polynomial algebra
$\C[p_1,p_2,\ldots]$.  The image under $\rho^R$ of any element of $\Lambda$
is therefore a bounded random variable.  This proves the required
$L^1$-specialization property.

Let $f\in\Lambda$.  After substituting \eqref{eq:centeredpk}, the random
variable $\rho^R(f)$ is an ordinary polynomial in $t_1,t_2$.  Applying
\eqref{eq:twomean} evaluates that polynomial at $t_1=t_2=c$.  At this point
every power sum in \eqref{eq:centeredpk} becomes zero, and hence
\begin{equation}\label{eq:counit}
 \E[\rho^R(f)]=f\big|_{p_1=p_2=\cdots=0}.
\end{equation}
For every nonempty partition $\lambda$, the Schur function $s_\lambda$ is
homogeneous of positive degree, so
\begin{equation}\label{eq:schurzero}
 \E[\rho^R(s_\lambda)]=0.
\end{equation}
All hook partitions in \eqref{eq:giambelli} are also nonempty.  Both sides of
\eqref{eq:giambelli} are consequently zero for every nonempty $\lambda$;
the empty partition gives $1=1$.  This proves Giambelli compatibility.

\subsection*{The iid two-point process is not determinantal}

Relative to the reference probability measure $\nu_R$, the ordered factorial
correlation densities of $X$ are
\begin{equation}\label{eq:corr}
 \rho_1(x)=2,\qquad \rho_2(x,y)=2\quad(x\ne y),
 \qquad \rho_m=0\quad(m\ge3).
\end{equation}
Indeed, the two labels give two choices for one ordered point and two choices
for an ordered pair, while no ordered tuple of three distinct points exists.

It suffices to rule out a kernel relative to $\nu_R$.  To see that this also
rules out every other reference measure, suppose that a kernel $K$ represents
the correlation measures relative to a measure $\mu$.  Since
$\rho_1=2\nu_R$, we have $\nu_R\ll\mu$; write $h=d\nu_R/d\mu$.  On the
$\nu_R$-conull set where $h>0$, the rescaled kernel
\begin{equation}\label{eq:rescale}
 B(x,y)=\frac{K(x,y)}{\sqrt{h(x)h(y)}}
\end{equation}
represents the same correlations relative to $\nu_R$, because every
$m\times m$ determinant is divided by $\prod_{j=1}^m h(x_j)$.

Assume for contradiction that such a measurable $B$ exists.  By
\eqref{eq:corr}, for almost every tuple of pairwise distinct points, its
principal minors satisfy
\begin{equation}\label{eq:minors}
 \det B_{\{j\}}=2,\qquad
 \det B_{\{j,k\}}=2,\qquad
 \det B_I=0\quad (|I|=3,4).
\end{equation}
There is a quadruple $(x_0,x_1,x_2,x_3)$ for which all the finitely many
relations in \eqref{eq:minors} hold simultaneously.  This follows by taking
a point in the intersection of the corresponding conull subsets of
$\R^4$; the atomlessness of $\nu_R$ removes the diagonals.  Write
$b_{jk}=B(x_j,x_k)$.

The one- and two-minor equations give
\begin{equation}\label{eq:pairs}
 b_{jj}=2,\qquad b_{jk}b_{kj}=2\quad(j\ne k).
\end{equation}
In particular every off-diagonal entry is nonzero.  Principal minors are
unchanged by diagonal similarity.  Conjugating the four-by-four matrix by a
diagonal matrix, we may therefore arrange
\begin{equation}\label{eq:gauge}
 b_{0j}=1,\qquad b_{j0}=2\quad(j=1,2,3).
\end{equation}

Fix distinct $j,k\in\{1,2,3\}$ and put $a=b_{jk}$.  Then
$b_{kj}=2/a$, and the zero determinant of the principal submatrix on
$\{0,j,k\}$ reads
\begin{equation}\label{eq:quadratic}
 0=\det\begin{pmatrix}2&1&1\\2&2&a\\2&2/a&2\end{pmatrix}
 =\frac{2(a^2-2a+2)}a.
\end{equation}
Hence $a\in\{1+i,1-i\}$.  Direct each edge $j\to k$ when
$b_{jk}=1+i$; by \eqref{eq:pairs}, the reverse entry is then $1-i$.
For the three vertices $1,2,3$, a directed three-cycle would give
\begin{equation}\label{eq:cyclic}
 \det\begin{pmatrix}
 2&1+i&1-i\\1-i&2&1+i\\1+i&1-i&2
 \end{pmatrix}=-8,
\end{equation}
contrary to the three-minor equation.  The tournament is therefore
transitive.  Relabeling $x_1,x_2,x_3$ if necessary, the full matrix is
\begin{equation}\label{eq:normalmatrix}
 M=\begin{pmatrix}
 2&1&1&1\\
 2&2&1+i&1+i\\
 2&1-i&2&1+i\\
 2&1-i&1-i&2
 \end{pmatrix}.
\end{equation}
A direct determinant calculation gives
\begin{equation}\label{eq:minusfour}
 \det M=-4.
\end{equation}
This contradicts the zero four-point correlation in
\eqref{eq:minors}.  No determinantal kernel exists, completing the proof of
Theorem~\ref{thm:main}. \qed

\section{Verification, scope, and novelty}

The accompanying exact-arithmetic script
\path{code/verify_counterexample.py} checks the determinant factorization in
\eqref{eq:quadratic}, the transitive and cyclic three-by-three determinants,
and $\det M=-4$.  It also constructs Schur functions from the Newton
recurrence and verifies the vanishing in \eqref{eq:schurzero} for every
partition through total degree eight.  The command used was
\begin{verbatim}
conda run --no-capture-output -n sandbox python \
  code/verify_counterexample.py
\end{verbatim}
and all exact checks passed.  These checks are not substitutes for the
all-degree mean-value proof or the matrix argument above.

The counterexample uses the continuous alternative $E=\R$ allowed by the
source.  It does not address whether a converse becomes true for a countable
discrete $E$, for uncentered positive specializations, or under additional
uniform-integrability or rigidity hypotheses.  Those strengthenings are
separate questions.

A bounded novelty search on 11 August 2026 checked the run indexes, the full
source of arXiv:2111.05606, the cited Borodin--Olshanski--Strahov source
arXiv:math-ph/0505021 \cite{bos}, and arXiv/web queries for the exact
conjecture and close variants.  No later claimed resolution or this Cauchy
counterexample was located.  The detailed search record is in
\path{novelty_search.md}; novelty confidence is moderate, not exhaustive.

\paragraph{Human-review recommendation.}
Review the circle pushforward and mean-value step, then check the normalization
under a change of reference measure and the four-point principal-minor
obstruction.  These are the only substantive mechanisms; no numerical or
unproved external input enters the counterexample.

\begin{thebibliography}{9}
\bibitem{source}
A.~I. Bufetov and P. Lazag,
\emph{A determinantal point process governed by an integrable projection
kernel is Giambelli compatible}, arXiv:2111.05606 (2021).

\bibitem{bos}
A. Borodin, G. Olshanski, and E. Strahov,
\emph{Giambelli compatible point processes},
Adv. Appl. Math. 37 (2006), 209--248; arXiv:math-ph/0505021.
\end{thebibliography}

\end{document}
